%% Mathematical appendix for the focused v82 revision.
\section{Mathematical foundations and proofs}\label{app:proofs}

\subsection{Conventions and conservation transfer}\label{app:notation}
The incidence matrix $B\in\R^{n\times m}$ has entries $-1$ at an arc's
origin and $+1$ at its destination; $\mu_j=e_{t_j}-e_{s_j}$ and
$M=[\mu_1,\ldots,\mu_K]$. Write $Z=\ker B$ and let $B^\dagger$ be the
Moore--Penrose pseudoinverse. All norms below are Euclidean or their
induced operator norms unless indicated otherwise. Graphs are loopless
multigraphs: parallel and reciprocal arcs remain distinct edges, including
in two-edge cycles. Cycle atoms and the block-cut tree are taken on the
underlying undirected multigraph. Formally, $\iota(v)$ is the cut-vertex
node when $v$ is an articulation vertex and otherwise its unique block
node; $\mathcal P(s,t)$ lists the atoms on the unique tree path from
$\iota(s)$ to $\iota(t)$. A signed simple $s$--$t$ path vector $p$ satisfies
$Bp=\mu$ and the full solution set is $p+Z$. Such a path uses only atoms
of $\mathcal P(s,t)$ and meets every atom on it.

\paragraph{Signed responses.}\label{app:signed}
The circulation and response vectors are signed first-order sensitivities,
not nonnegative commodity flows. At a strictly positive flow, sufficiently
small signed circulation perturbations also preserve nonnegativity. A
restricted response family is treated explicitly below.

\begin{corollary}[IG exchange]\label{cor:ig}
For an exactly conservative $C^1$ field on an open neighborhood of the
attribution path, $BG_{:,j}=\mu_j\Delta q_j$.
\end{corollary}
\begin{pf}
Differentiate $Bf(q)=Mq$ to obtain $BJ_f(q)=M$, then integrate the
$j$th column along $\gamma(\lambda)=q^0+\lambda\Delta q$.
\end{pf}

\subsection{Cycle structure and the classification proofs}\label{app:classificationproof}
\begin{lemma}[structure of circuit classes]\label{lem:mprime}
The non-bridge atoms are the non-singleton circuit classes; each bridge
is a singleton class. Moreover $Z=\bigoplus_A Z_A$, where
$Z_A=\{z\in Z:\operatorname{supp}z\subseteq A\}$, and the summands
have disjoint edge supports.
\end{lemma}
\begin{pf}
No simple cycle crosses an articulation vertex into another block: it
would revisit that vertex. Within a non-bridge block, any two edges lie
on a common cycle. One way to see this is to subdivide the two edges;
the resulting block is still biconnected, so its two new vertices lie
on a common cycle, which descends to the required cycle. This is the
usual circuit-class characterization of the graphic matroid
\citep{oxley2011matroid}. Every signed circulation is a linear combination
of signed simple-cycle vectors: subtract the fundamental cycles of a
spanning forest to leave a circulation on a forest, necessarily zero
by successive leaf elimination. Grouping these cycles by their blocks
gives the direct sum, and distinct blocks have disjoint edge sets.
\end{pf}

\begin{lemma}[circulation-preserving diagonal maps]\label{lem:m}
For $D=\operatorname{diag}(w)$, $DZ\subseteq Z$ if and only if $w$ is
constant on every non-bridge atom; bridge weights are free.
\end{lemma}
\begin{pf}
For any signed simple-cycle vector $z_C$, a circulation supported on
that cycle has equal signed amplitudes on successive edges, by node
balance. Thus $BDz_C=0$ forces $w$ constant on $C$, and
Lemma~\ref{lem:mprime} propagates equality throughout each atom.
Conversely, atomwise constancy makes $D$ a scalar on every $Z_A$, so
it preserves their direct sum. Bridges carry no circulation.
\end{pf}

\begin{pf}[Theorem~\ref{thm:t3cpmain}]
Add a virtual reverse OD edge $(t,s)$ and set
\[
\bar B=[B,-\mu],\qquad \bar D=\operatorname{diag}(w,\alpha).
\]
For a fixed $\alpha$, property~\eqref{eq:t3prop} is equivalent to
$\bar D\ker\bar B\subseteq\ker\bar B$.
Indeed, a kernel vector $(x,a)$ obeys $Bx=a\mu$. If $a\ne0$,
apply \eqref{eq:t3prop} to $x/a$ to obtain $BDx=a\alpha\mu$.
If $a=0$, subtract the identities for $p+x$ and $p$, with any fixed
$Bp=\mu$, to obtain $BDx=0$. These are exactly the two cases of
$\bar B\bar D(x,a)=0$. Conversely, applying kernel invariance to
$(x,1)$ gives \eqref{eq:t3prop} for every $Bx=\mu$, at that same $\alpha$.

By Lemma~\ref{lem:m}, this invariance holds precisely when the augmented
weights are constant on its cycle atoms. The virtual edge and a simple
$s$--$t$ path form a cycle meeting every original atom on
$\mathcal P(s,t)$, so these atoms merge with the virtual edge into one
circuit class. No other atoms join it: every new simple cycle contains
the virtual edge and a simple $s$--$t$ path, which uses only those
atoms. Other original atoms retain their classes. Thus the augmented
criterion is exactly within-atom constancy everywhere and common value
$\alpha$ on $\bigcup\mathcal P(s,t)$. This proves both directions and
identifies $\alpha=w^*$ without changing the universal response family.
\end{pf}

\begin{pf}[Theorem~\ref{thm:t3cm} and Corollary~\ref{cor:path}]
Apply the single-OD theorem to each query at the same $w$. Necessity
requires equality within every atom and among all atoms co-traversed
by a queried block path. Their transitive closure is exactly the
compatibility-class partition. Conversely, class constancy supplies
both single-OD conditions for every query. Along the attribution
path, exact conservation gives $BJ_f(\gamma(\lambda))\mathbf e_j=\mu_j$;
the static theorem therefore yields
\[
B\widetilde C_{:,j}=\Delta q_j\int_0^1
BD_\lambda J_f(\gamma(\lambda))\mathbf e_j\,d\lambda
=\Delta q_j\left(\int_0^1w^*_{k(j)}(\lambda)\,d\lambda\right)\mu_j.
\]
\end{pf}

\begin{corollary}[all-OD class collapse]\label{cor:allod}
Requiring the criterion at one common $w$ for every ordered pair of
distinct vertices makes all atoms in a connected component share one
constant.
\end{corollary}
\begin{pf}
For two atoms incident to the same cut vertex, choose endpoints on
their opposite sides. Their OD block path contains both atoms, forcing
their constants to agree. Connectivity of the block-cut tree propagates
the equality throughout the component.
\end{pf}

\begin{rmk}[a common flow realization]\label{rmk:quantifier}
The merging proof requires one common weight vector. Computing the
queries at different realized flows does not in general supply this
premise.
\end{rmk}

\begin{table}[pos=tb]
\caption{Constructive query-dependent compatibility on Figure~\ref{fig:blockcut}
of the main text. Ranges use all $512$ separating sets per query and a
conservative unit response; the weighted circulation residual is below
$5\times10^{-16}$ in all stages, so these ranges are independent of that
response. The intentionally chosen linear weights illustrate
compatibility; they are not the marginal costs of the original TSTT
outcome.}\label{tab:topologyexample}
\begin{tabular*}{\tblwidth}{@{\extracolsep{\fill}}lccccl}
\toprule Stage & $w_{A,b_1,B}$ & $w_{b_2,C}$ & $w_D$ & Queried ODs & Account ranges\\\midrule
Initial query & 1 & 2 & 3 & $(s,u)$ & $[1,1]$\\
Add query & 1 & 2 & 3 & $(s,u);(u,t)$ & $[1,1];[1,2]$\\
Adjust merged class & 1 & 1 & 3 & $(s,u);(u,t)$ & $[1,1];[1,1]$\\
\bottomrule\end{tabular*}\end{table}
The calculation uses the minimum-norm solution of $Bx=\mu$ and checks
$BDZ=0$ on a basis of $\ker B$, which establishes response independence
of the displayed boundary vector. The theorem supplies the uniform
common-account guarantee in stages one and three.

\subsection{Direction, integration, and network representation}\label{supp:criteria}
At fixed weights $D=\operatorname{diag}(w)$, a realized model balances
when $BDJ_j=\alpha_j\mu_j$ for its actual response.
\emph{Universal directional balance} requires
$BDx\in\operatorname{span}\{\mu_j\}$ for every $Bx=\mu_j$, allowing its
coefficient to depend on $x$; \emph{structural cut consistency}
additionally requires one common coefficient across all such responses,
and is the property classified in the main text. For non-collinear
queries with shared weights, universal directional balance already
implies structural cut consistency because
$BD\ker B\subseteq\bigcap_j\operatorname{span}\{\mu_j\}=\{0\}$;
a query and its reverse alone do not suffice.
\begin{proposition}[universal directional criterion]\label{prop:directional}
Let $\hat x_j=B^\dagger\mu_j$, let $N$ span $Z$ with orthonormal
columns, and let $P_{\mu_j}$ project orthogonally onto
$\operatorname{span}\{\mu_j\}$. At fixed $w$, universal directional
balance is equivalent to
\[
(I-P_{\mu_j})BD\hat x_j=0,\qquad (I-P_{\mu_j})BDN=0.
\]
The scale may depend on the response. Requiring one common boundary value
additionally requires $BDN=0$. If directional balance holds along the
whole path for a conservative model, its integrated column balances.
\end{proposition}
\begin{pf}
Every response is $\hat x_j+z$, $z\in Z$. Apply $I-P_{\mu_j}$ to
$BD(\hat x_j+z)$ and vary $z$. The scale is
$\mu_j^\top BD(\hat x_j+z)/\norm{\mu_j}^2$ and is independent of
$z$ precisely when $BDZ=0$. Integrating the realized directional
identity gives the endpoint conclusion. Equivalently, invariance on
all incidence fibers, $Bx=By\Rightarrow BDx=BDy$, is exactly
$DZ\subseteq Z$, since $x-y$ ranges over $Z$.
\end{pf}

\begin{example}[strict separation of the criteria]\label{ex:triangle}
On the diamond $(s,a),(a,t),(s,b),(b,t)$ with weights $(1,1,1,2)$,
the upper-path response $(1,1,0,0)$ has divergence $\mu$ after
weighting, but the lower-path response has weighted divergence
$(-1,0,-1,2)$ in node order $(s,a,b,t)$, outside $\operatorname{span}\{\mu\}$.
On the triangle $(s,t),(t,u),(u,s)$ with weights $(1,0,0)$, every
response $x=(1+c,c,c)$ satisfies $BDx=(1+c)\mu$; direction balances,
but the common boundary value fails because the circulation $(1,1,1)$ maps
to $\mu\ne0$.
\end{example}

If directional balance holds for all
queries at common $w$, differences give
$BDZ\subseteq\bigcap_j\operatorname{span}\{\mu_j\}$. For non-collinear
queries this intersection is $\{0\}$, so directional balance implies
the common-value property. This uses the full response space $Z$.

\paragraph{Restricted families.}\label{app:membership}
For $S\subseteq Z$ with basis $N_S$, test responses in
$x_j^{\mathrm{ref}}+S$, where $Bx_j^{\mathrm{ref}}=\mu_j$.
The same proof gives the direction tests with $\hat x_j,N$ replaced
by $x_j^{\mathrm{ref}},N_S$; common boundary value replaces the second test by
$BDN_S=0$. Using $\hat x_j$ requires
$\hat x_j-x_j^{\mathrm{ref}}\in S$. For a fixed-path model $f=\Delta h$,
$A_{\mathrm{OD}}h(q)=q$ and $B\Delta=MA_{\mathrm{OD}}$. Choose a path
vector $a_j$ with $A_{\mathrm{OD}}a_j=\mathbf e_j$ and set
$x_j^{\mathrm{ref}}=\Delta a_j$. Differentiation yields
\[
J_f\mathbf e_j-x_j^{\mathrm{ref}}
=\Delta(\partial_jh-a_j)\in\Delta\ker A_{\mathrm{OD}}=:T_{\mathrm{swap}}.
\]
This space includes OD path exchanges of every OD: background
ODs may reroute even when their demands remain fixed.

\begin{example}[endpoint cancellation]\label{ex:cancellation}
For two series bridges, $f(q)=(q,q)$ and
$T(F)=F_1+F_2+[\sin(2\pi F_1)-\sin(2\pi F_2)]/(2\pi)$ give
\[
\widetilde C_1-\widetilde C_2
=\frac{\sin(2\pi q^1)-\sin(2\pi q^0)}{\pi}.
\]
For $q^0=0$, the nonzero balanced endpoints in $[0,1]$ are
$q^1=1/2,1$, although the two marginal weights differ almost
everywhere. At $q^1=0.3$ balance fails. Endpoint cancellation therefore
does not establish the pointwise classification conditions.
\end{example}

\begin{example}[network subdivision]\label{ex:subdivision}
A single edge with $f(q)=q$, $T(F)=\tau F$, $\tau>0$, has cut account
$\tau\Delta q$. Subdivide it into $(s,v),(v,t)$ and set
$f'(q)=(q,q)$,
$T'(F_1,F_2)=\tau[\alpha F_1+(1-\alpha)F_2]$, $0<\alpha<1$.
The composite demand-to-outcome map is unchanged, but
$\widetilde C'=\tau\Delta q(\alpha,1-\alpha)$ and
$(B\widetilde C')_v=\tau(2\alpha-1)\Delta q$.
The two cuts account $\tau\alpha\Delta q$ and
$\tau(1-\alpha)\Delta q$. They agree only for $\alpha=1/2$, when each
is half the undivided account; their sum remains the scalar
contribution $\tau\Delta q$. The invariance studied here thus concerns
conservative responses on a fixed network with a fixed edgewise
outcome definition, rather than invariance under subdivision.
\end{example}

\subsection{Approximate conservation and attribution balance}\label{app:approx}\label{app:failurestatement}
The residual identity~\eqref{eq:residualidentity} shows that output residuals alone do not control attribution columns. For unit $v\in\operatorname{im}B$ and
$r_\varepsilon(q)=\varepsilon v\tanh((q_1-q_2)/\varepsilon)$, the field $f_\varepsilon(q)=B^\dagger(Mq+r_\varepsilon(q))$ has residual exactly $r_\varepsilon$, globally bounded by $\varepsilon$ and zero on $q_1=q_2$; yet on $q(\lambda)=\lambda(\delta,\delta)$ its partial derivatives are $v,-v$, so $R_1=\delta v$ and $R_2=-\delta v$, and its residual Jacobian Lipschitz constant grows as $\varepsilon^{-1}$. Added to an exactly conservative field with minimum flow $\varrho>0$ on a compact positive-demand box, it preserves strictly positive flows when $\varepsilon\|B^\dagger\|_2<\varrho$. The traffic experiments use exactly conservative models; these bounds delimit scope rather than enter the numerics.

\begin{theorem}[local neighborhood bound]\label{thm:t1pmain}
Let $r\in C^{1,1}(\mathcal N;\R^n)$ on an open neighborhood of the attribution path, with $\operatorname{Lip}(J_r)\le H$. Suppose, for some $\rho>0$ and every $j,\lambda$, the coordinate segment $\gamma(\lambda)+[-\rho,\rho]\mathbf e_j$ lies in $\mathcal N$ and its two endpoint residual norms are at most $\varepsilon$. Then
\[
\|R_j\|\le|\Delta q_j|(\varepsilon/\rho+H\rho/2).
\]
\end{theorem}
\begin{pf}
At $q=\gamma(\lambda)$, Taylor's formula gives $r(q\pm\rho\mathbf e_j)=r(q)\pm\rho\partial_jr(q)+\xi_\pm$ with $\|\xi_\pm\|\le H\rho^2/2$. Subtraction bounds $\|\partial_jr(q)\|$ by $\varepsilon/\rho+H\rho/2$; integration gives the claim.
\end{pf}
The optimized square-root rate is sharp. Indeed $r_\varepsilon=\varepsilon v\sin((q_1-q_2)/\sqrt\varepsilon)$ has supremum norm $\varepsilon$ and $\operatorname{Lip}(J_{r_\varepsilon})\le2$, while the same diagonal path gives $R_1=\delta\sqrt\varepsilon v=-R_2$. Thus a uniform $o(\sqrt\varepsilon)$ rate is impossible for bounded Jacobian Lipschitz constant.

\begin{lemma}[minimum-norm support]\label{lem:s}
For a signed $s$--$t$ path $p$, $\hat x=B^\dagger\mu=(I-P_Z)p$
is supported on $\bigcup\mathcal P(s,t)$.
\end{lemma}
\begin{pf}
The minimum-norm point of $p+Z$ is its orthogonal projection onto
$Z^\perp$. Lemma~\ref{lem:mprime} gives
$P_Z=\sum_A P_{Z_A}$ with disjoint supports; both $p$ and every
$P_{Z_A}p$ vanish on off-path atoms.
\end{pf}

\begin{proposition}[combined dispersion--residual bound]\label{prop:combined}
Let $f\in C^1$ near the attribution path, $r=Bf-Mq-b$ with the constant
background balance $b$ of Section~\ref{sec:noninheritance} ($b=0$
without background), and let
$w(\lambda)=\nabla T(f(\gamma(\lambda)))$ be continuous.
Let $\bar w(\lambda)$ be the compatibility-class mean vector and
$\eta(\lambda)=\norm{w(\lambda)-\bar w(\lambda)}_\infty$.
Writing $J_j(\lambda)=J_f(\gamma(\lambda))\mathbf e_j$ and
$d_j=(I-P_{\mu_j})B\widetilde C_{:,j}$, we have
\[
\norm{d_j}_2\le |\Delta q_j|\norm{B}_2\int_0^1\left[
\eta(\lambda)\norm{J_j(\lambda)}_2+
\norm{\bar w(\lambda)}_\infty\norm{B^\dagger}_2
\norm{\partial_jr(\gamma(\lambda))}_2\right]d\lambda.
\]
\end{proposition}
\begin{pf}
Differentiation gives $BJ_j=\mu_j+\partial_jr$ --- $b$ is constant and
drops out of the derivative --- so
$J_j=\hat x_j+z_j+B^\dagger\partial_jr$, $z_j\in Z$.
Let $\bar D=\operatorname{diag}(\bar w)$. Class constancy gives
$B\bar Dz_j=0$ by Lemma~\ref{lem:m}. By Lemma~\ref{lem:s},
$\hat x_j$ is supported on the OD block path, whose atoms lie in one
compatibility class $k(j)$; hence
$B\bar D\hat x_j=\bar w_{k(j)}\mu_j$.
Therefore
\[
(I-P_{\mu_j})BDJ_j=(I-P_{\mu_j})B
\left[\operatorname{diag}(w-\bar w)J_j+
\bar DB^\dagger\partial_jr\right].
\]
The projector is non-expansive, the two diagonal operator norms are
$\eta$ and $\norm{\bar w}_\infty$, and integration gives the bound.
For exact conservation only the dispersion term remains; for
$\eta=0$ only the derivative residual remains. Theorem~\ref{thm:t1pmain}
can bound that derivative term under its local hypotheses.
\end{pf}

\subsection{Exact boundary error and its path bound}\label{app:boundaryerror}
\begin{pf}
Set $b=Be_j$ and $I=V\setminus\{s_j,t_j\}$. Since $\mathbf1^\top b=0$,
Eq.~\eqref{eq:cutextrema} applied to $e_j$ gives an interval of extrema
with midpoint
\[
-b_{s_j}-\frac12\sum_{v\in I}b_v
=\frac12(b_{t_j}-b_{s_j})
\]
and half-width $\frac12\sum_{v\in I}|b_v|$. The larger absolute
endpoint is the absolute midpoint plus this half-width, proving
Eq.~\eqref{eq:exactboundaryerror}. This argument concerns the attained
extrema; it does not require every value between them to be attained.
The map $\mathcal K_j$ is a seminorm because it is a sum of absolute
linear functionals. Moreover
$e_j=\Delta q_j\int_0^1 h_j(\gamma(\lambda))\,d\lambda$.
Its homogeneity and triangle inequality therefore yield
\[
\epsilon_j^{\mathcal U}\le\mathcal K_j(e_j)
\le |\Delta q_j|\int_0^1\mathcal K_j(h_j(\gamma(\lambda)))\,d\lambda,
\]
proving the path statement. Its kernel is $\ker B$: if the seminorm
vanishes, all interior divergences vanish and $b_{s_j}=b_{t_j}$;
the zero sum then makes both endpoints zero.
\end{pf}

